\section{The multiscale potential and the canonical singularity}
\label{sec:iteration}

Fix \(t\) as in \Cref{i:initial-separation}.  We will preserve its
three groups of generating supports while replacing ideal attractions
by physical ones.  The essential step is quantitative: a pressure error
\(\delta\), followed by an ideal attraction of mass \(a\), enlarges
the slope groups by \(O(\delta/a+a)\) and costs \(O(a)\) in the
parameter domain.  We first prove this step, then choose one sequence
of masses and ranges for which both losses are summable.

\subsection{Preservation under one attraction}

Let \(q\) be a physical pressure and \(P\) a convex approximation
whose subdifferential is generated near the three initial slopes.
In the iteration, \(P\) will be the preceding ideal transform and
\(q\) its finite-range realization.  We ask how closely the generating
supports of \(T_aq\) remain in those three groups when
\(\norm{P-q}_{W_1}\le\delta\).  At each active shifted point, the
parabolic bound first isolates one group in a neighborhood of radius
proportional to \(a\).  Secant inequalities in that neighborhood then
transfer the support bound from \(P\) to \(q\).

Recall the parameter balls
\(W(r)=\clB{\lambda_*}{r}\), \(W_1=W(1/4)\),
the slope vectors \(v_i=(Ad_i^2/2,d_i)\), and the constants
\(\tau,\gamma>0\) from \Cref{i:geometry}.
By \Cref{k:transform-properties}, choose \(U_*>0\) bounding all active
shifts for every transform \(T_aQ\) of a family pressure, with \(0<a\le A+1\), at parameters
in \(W_1\), uniformly over admissible references and masses.
Uniform convex-function bounds on a larger compact neighborhood give a
common Lipschitz constant \(C_*\ge1\) for all these transforms on \(W_1\).
These constants depend only on the fixed family and parameter compacts.
Put \(\beta_{\min}=3/4\), \(\beta_{\max}=5/4\).
Choose \(k>0\), and then \(l>0\), sufficiently small that
\begin{equation}\label{eq:band-constants}
 \frac{k}{2\beta_{\min}}+\frac{2C_*l}{k}<\frac p{16}.
\end{equation}

Set
\[
 S_* = \beta_{\max}U_*+k+l+1.
\]
Here \(\beta_{\max}U_*a\) bounds the displacement to an active
shifted point, \(ka\) will be a field increment, and \(la\) the
radius on which all supports will be controlled.  All these constants
depend only on the fixed \(t\), the family, and the parameter compacts.

\begin{lemma}[One-step preservation of generating supports]
\label{m:one-step}
Let \(q\) be a pressure from the family \eqref{f:family}, and let
\(P\) be a finite convex function on \(\Omega\) that is
\(C_*\)-Lipschitz on \(W_1\).  Suppose
\[
 \begin{gathered}
 0<a\le A+1,\qquad 0<r<\tau,\qquad S_*a<s\le\frac14,\\
 \delta\ge0,\qquad \norm{P-q}_{W_1}\le\delta,\qquad
 \frac{2\delta}{ka}<\frac p{16}.
 \end{gathered}
\]
Assume that throughout \(W(s)\),
\begin{equation}\label{eq:one-step-hypothesis}
 \partial P(y)=\conv\left(\partial P(y)\cap
                 \bigcup_{i=1}^3\clB{v_i}{r}\right).
\end{equation}
Then, with
\[
 r^+=r+\frac{2\delta}{la}+\frac12aU_*^2,
\]
one has for every \(\lambda\in W(s-S_*a)\),
\begin{equation}\label{eq:one-step-conclusion}
 \partial(T_aq)(\lambda)
 =\conv\left(\partial(T_aq)(\lambda)\cap
                 \bigcup_{i=1}^3\clB{v_i}{r^+}\right).
\end{equation}
\end{lemma}

\begin{proof}
Fix \(\lambda=(\beta,h)\in W(s-S_*a)\), and choose an active
shift \(u\) for \(T_aq\) at \(\lambda\).
Let \(x=(\beta,h+\beta au)\).
The distance from \(\lambda\) to the boundary of \(W(s)\) is at least
\(S_*a\).
Since \(\abs{x-\lambda}\le\beta_{\max}U_*a\), every point
\[
 y,\quad y+(0,\pm ka),\quad x+(0,\pm ka),
 \qquad \abs{y-x}\le la,
\]
lies in \(W(s)\).
The surplus \(a\) leaves these test points in the interior of \(W(s)\).

\emph{A common density band.}
For \((e,d)\in\partial P(y)\), the supporting inequality at \(y\),
the Lipschitz bound on \(P\), and
the bound \(\norm{P-q}_{W_1}\le\delta\) give
\begin{align*}
 \pm ka\,d
 &\le P(y+(0,\pm ka))-P(y)\\
 &\le q(x+(0,\pm ka))-q(x)
                         +2\delta+2C_*la\\
 &\le \pm ka\,u+\frac{(ka)^2}{2\beta a}
                         +2\delta+2C_*la.
\end{align*}
The last step is the active-point parabola of \Cref{k:active-parabola}.
Using both signs, \eqref{eq:band-constants}, and the hypothesis on
\(\delta/a\), we obtain
\begin{equation}\label{eq:common-density-band}
 \abs{d-u}\le
 \frac{k}{2\beta_{\min}}+\frac{2C_*l}{k}
                       +\frac{2\delta}{ka}<\frac p8.
\end{equation}
The same \(u\) works for every \(y\) in this neighborhood.

\emph{One slope group near the shifted point.}
The set in the right side of \eqref{eq:one-step-hypothesis} is nonempty:
the full subdifferential is nonempty, and the assumed hull identity
identifies it with the convex hull of that set.
If a generator belongs to \(\clB{v_i}{r}\), then
\[
 \abs{u-d_i}<\frac p8+\tau<\frac{3p}{16}.
\]
Since the distinct \(d_i\) are separated by at least \(p\), at most one
index \(i\) can satisfy this condition.
It is therefore one common index for the entire neighborhood.
By convexity of the ball and the assumed hull identity,
\begin{equation}\label{eq:one-ball-neighborhood}
 \partial P(y)\subset\clB{v_i}{r}
 \qquad(\abs{y-x}\le la).
\end{equation}

\emph{Transfer to the physical pressure and its transform.}
Take \(v\in\partial q(x)\) and a unit vector \(z\in\R^2\).
At \(y=x+la z\), choose any support of \(P\).
Using \eqref{eq:one-ball-neighborhood} to bound this support and
the approximation at \(x,y\), we get
\[
 la\,v\cdot z
 \le q(y)-q(x)
 \le P(y)-P(x)+2\delta
 \le la\,(v_i\cdot z+r)+2\delta.
\]
Testing every unit \(z\) proves
\[
 \abs{v-v_i}\le r+\frac{2\delta}{la}.
\]
By \Cref{k:subdifferential}, the generating supports of the active branch
of \(T_aq\) are obtained by adding \((au^2/2,0)\).
They all lie in \(\clB{v_i}{r^+}\) by
the definition of \(r^+\) and \(\abs{u}\le U_*\).
The exact convex-hull formula over all active branches gives
\eqref{eq:one-step-conclusion}.
\end{proof}

The common density band is what prevents supports near one active
shifted point from mixing different groups.  Once one group has been
selected, ordinary secants transfer its entire support bound across a
pressure error \(\delta\); division by the test distance \(la\)
explains the term \(2\delta/(la)\).  The remaining term in \(r^+\)
bounds the shift \((au^2/2,0)\) of each support under the transform.

\subsection{A common choice of scales}\label{m:schedule}

For \(R_1>2\), define
\begin{equation}\label{eq:mass-range}
 R_n=R_1^{2^{n-1}},\qquad
 a_1=A,\qquad a_n=R_n^{-1/16}\quad(n\ge2).
\end{equation}
Let
\begin{equation}\label{eq:truncated-models}
 \varphi_n=\psi_t-t\sum_{j=1}^n a_jw_{R_j},
 \qquad q_n=Q_{\varphi_n},
 \qquad P_n=T_{a_n}q_{n-1}\quad(n\ge1),
\end{equation}
where \(\varphi_0=\psi_t\) and \(q_0=q^{(t)}\).
Thus \(P_1=T_Aq^{(t)}\) is the pressure used in
\Cref{i:initial-separation}.

Choose \(R_1\) large enough initially that
\(\sum_{n\ge2}a_n<1\).
Every truncated potential then belongs to the stable family
\eqref{f:family}.
By \Cref{k:range}, for all sufficiently large \(R_1\),
\begin{equation}\label{eq:approximation-errors}
 \norm{q_n-P_n}_{W_1}\le\delta_n,
 \qquad \delta_n=C_3R_n^{-1/4}\quad(n\ge1),
\end{equation}
with one constant \(C_3\) independent of \(n\) and \(R_1\).
Indeed, the reference range is at most two for \(n=1\), giving the
stronger estimate \(O(R_1^{-1/2})\).
For \(n\ge2\), it is at most \(R_{n-1}=\sqrt{R_n}\).
The ratio threshold in \Cref{k:range} is uniform in the reference, so it
suffices to enlarge the same \(R_1\).

The crucial relations are
\begin{equation}\label{eq:summable-errors}
 a_{n+1}=R_n^{-1/8},\qquad
 \frac{\delta_n}{a_{n+1}}=C_3R_n^{-1/8},
 \qquad
 \sum_{n\ge1}\left(a_{n+1}+\delta_n+
             \frac{\delta_n}{a_{n+1}}\right)
       \longrightarrow0\quad(R_1\to\infty).
\end{equation}
For example, for every \(q>0\),
\[
 \sum_{n\ge1}R_n^{-q}
 =\sum_{n\ge1}(R_1^{-q})^{2^{n-1}}
 \le \frac{R_1^{-q}}{1-R_1^{-q}}.
\]
This explicit bound will meet all the conditions below.

By \Cref{k:transform-properties}, fix \(C_0>0\), independently of
\(R_1\), such that
\begin{equation}\label{eq:small-mass-increment}
 0\le T_aQ-Q\le C_0a\qquad\text{on }W_1
\end{equation}
for all family pressures and \(0<a\le A+1\).

Set
\[
 W_n=W\left(\frac14-S_*\sum_{j=1}^{n-1}a_{j+1}\right),
\]
and define
\begin{equation}\label{eq:slope-radii}
 r_1=\frac{\tau}{2},\qquad
 r_{n+1}=r_n+\frac{2\delta_n}{la_{n+1}}
                   +\frac12a_{n+1}U_*^2.
\end{equation}

We now enlarge \(R_1\) once more so that
\begin{align}
 S_*\sum_{n\ge1}a_{n+1}&<\frac18,
 &\sum_{n\ge1}\left(\frac{2\delta_n}{la_{n+1}}
                         +\frac12a_{n+1}U_*^2\right)&<\frac{\tau}{2},
 \label{eq:domain-radius-budget}\\
 \sup_{n\ge1}\frac{2\delta_n}{ka_{n+1}}&<\frac p{16},
 &\sum_{n\ge1}(\delta_n+C_0a_{n+1})&<\frac{\gamma}{64}.
 \label{eq:band-pressure-budget}
\end{align}
All four conditions follow from \eqref{eq:summable-errors};
no further choices of masses or ranges are made.
In particular, every \(W_n\) contains \(W(1/8)\) and \(r_n<\tau\).

The pressures stay uniformly close to the three-plane limit:
\begin{equation}\label{eq:pressure-near-limit}
 \norm{P_n-P_*}_{W_1}<\frac{\gamma}{32}\qquad(n\ge1).
\end{equation}
To see this, use
\[
 \norm{P_{n+1}-P_n}_{W_1}
 \le \norm{T_{a_{n+1}}q_n-q_n}_{W_1}
          +\norm{q_n-P_n}_{W_1}
 \le C_0a_{n+1}+\delta_n
\]
and telescope from the initial error \(\gamma/64\).
This comparison uses \eqref{eq:small-mass-increment} for the physical
pressure \(q_n\), so no shifted evaluation of the approximation
\eqref{eq:approximation-errors} is required here.

\begin{proposition}[Persistence at every scale]\label{m:propagation}
For every \(n\ge1\) and \(y\in W_n\),
\begin{equation}\label{eq:separated-generators}
 \partial P_n(y)=
 \conv\left(\partial P_n(y)\cap
                 \bigcup_{i=1}^3\clB{v_i}{r_n}\right).
\end{equation}
\end{proposition}

\begin{proof}
The case \(n=1\) is \Cref{i:initial-separation}.
Each \(q_n\) belongs to the family, and each \(P_n\) has the common
Lipschitz bound \(C_*\).  Apply \Cref{m:one-step} inductively with
\(P=P_n\), \(q=q_n\), \(a=a_{n+1}\), \(\delta=\delta_n\),
\(r=r_n\), and \(s\) equal to the radius of \(W_n\).
The first condition in \eqref{eq:domain-radius-budget} keeps the
parameter radius above \(1/8\), while the second gives \(r_n<\tau\).
The other hypotheses are \eqref{eq:approximation-errors} and the first
condition in \eqref{eq:band-pressure-budget}.  The new domain and
slope radius are exactly \(W_{n+1}\) and \(r_{n+1}\).
\end{proof}

The three groups now force a two-dimensional set of supports at one
parameter.  Minimizing against their barycenter first finds that
barycenter in a subdifferential; its separation from every two-group
hull then forces all three groups to occur there.

\begin{proposition}[A common ball of pressure supports]\label{m:pressure-ball}
For every \(n\ge1\), there is
\(\lambda_n\in\operatorname{int}W(1/8)\) such that
\begin{equation}\label{eq:common-pressure-ball}
 \clB{g}{\gamma}\subset\partial P_n(\lambda_n).
\end{equation}
\end{proposition}

\begin{proof}
Minimize \(P_n(\lambda)-g\cdot\lambda\) on the compact ball \(W(1/8)\).
For a boundary point, \Cref{i:geometry} and
\eqref{eq:pressure-near-limit} give
\begin{align*}
 &[P_n(\lambda)-g\cdot\lambda]
   -[P_n(\lambda_*)-g\cdot\lambda_*]\\
 &\hspace{2em}>
   \frac{\gamma}{8}-\frac{2\gamma}{32}
   =\frac{\gamma}{16}>0.
\end{align*}
Thus a minimizer \(\lambda_n\) is interior, and
\(g\in\partial P_n(\lambda_n)\).
By \Cref{m:propagation}, this subdifferential is generated by its
intersections with the three balls \(\clB{v_i}{r_n}\).
The vector \(g\) lies outside the convex hull of any two of the larger
balls \(\clB{v_i}{\tau}\).
Each of the three intersections must therefore be nonempty.
Choose one actual support from each.
Their triangle contains \(\clB{g}{\gamma}\) by \Cref{i:geometry}, and
their convex hull is contained in \(\partial P_n(\lambda_n)\).
\end{proof}

\subsection{The final potential and its decay}

Define one potential, now with all parameters fixed, by
\begin{equation}\label{eq:final-potential}
 \Phi(x)=\phi(\abs{x})
   =\psi_t(\abs{x})-t\sum_{n\ge1}a_nw_{R_n}(x).
\end{equation}

\begin{proposition}[Stability and algebraic decay]\label{m:tail}
The function \(\phi\) is bounded and continuous on \([0,\infty)\),
is stable with \(B=t^2/2\), and satisfies
\(\abs{\phi(r)}\le Cr^{-3-1/32}\) for all \(r\ge1\), where
\begin{equation}\label{eq:tail-constant}
 C=2^{3+1/32}t^2+t\norm{w}_{\infty}
       \left(A R_1^{1/32}
              +\sum_{n\ge2}R_1^{-2^{n-1}/32}\right)<\infty.
\end{equation}
\end{proposition}

\begin{proof}
Since \(\sum_na_n<A+1\) and \(R_n\ge1\), the attraction series
converges uniformly.
The core and every summand are bounded, continuous, and radial.
The common cell estimate of \Cref{f:stability} applies to the sum and
gives \(U^\Phi\ge-B_0N=-(t^2/2)N\).

Set \(\varepsilon=1/32\).
A nonzero summand at radius \(r\) must have \(R_n\ge r\), and then
\[
 r^{3+\varepsilon}R_n^{-3}\le R_n^\varepsilon.
\]
The core vanishes for \(r\ge1+1/t<2\) and is at most \(t^2\).
For every \(r\ge1\), therefore,
\[
 r^{3+\varepsilon}\abs{\phi(r)}
 \le 2^{3+\varepsilon}t^2
          +t\norm{w}_\infty\sum_{n:R_n\ge r}a_nR_n^\varepsilon
 \le C.
\]
Finally,
\[
 \sum_{n\ge1}a_nR_n^{1/32}
 =A R_1^{1/32}+\sum_{n\ge2}R_1^{-2^{n-1}/32}<\infty
\]
by the geometric-series bound following
\eqref{eq:summable-errors}.
\end{proof}

\subsection{From pressure supports to a temperature singularity}

By \Cref{f:thermodynamics}, the infinite potential has canonical and
grand-canonical thermodynamic limits.
By \Cref{f:tail-pressure} and \eqref{eq:approximation-errors},
\[
 \norm{P_n-Q_\Phi}_{W_1}
 \le \delta_n+
 C_{W_1,t}\sum_{j>n}a_j\longrightarrow0.
\]
Choose a subsequence of \(\lambda_n\) converging to
\(\lambda_c=(\beta_c,h_c)\in W(1/8)\).
In particular, \(\beta_c\in[7/8,9/8]\).
For each \(v\in\clB{g}{\gamma}\) and \(\lambda\in W_1\),
\eqref{eq:common-pressure-ball} gives
\[
 P_n(\lambda)\ge P_n(\lambda_n)+v\cdot(\lambda-\lambda_n).
\]
Uniform convergence and continuity pass this inequality to the limit.
It is a supporting inequality for \(Q_\Phi\) on a neighborhood of
\(\lambda_c\).
Such a local support extends globally on the convex domain \(\Omega\):
a violation at another point would, by convexity along the joining
segment, give a violation arbitrarily close to \(\lambda_c\).
Consequently,
\begin{equation}\label{eq:limiting-support-ball}
 \clB{g}{\gamma}\subset\partial Q_\Phi(\beta_c,h_c).
\end{equation}

We state explicitly the duality implication used to finish the proof.
It is the supporting-hyperplane form of Fenchel duality
\citep[Theorem~11.1 and Proposition~11.3]{RockafellarWets1998}, applied
only in the field coordinate and with the inverse temperature retained
as a parameter.

\begin{lemma}[From pressure supports to canonical slopes]\label{m:canonical-corner}
Suppose that \(H(\beta,d)\) and \(Q(\beta,h)\) have the duality and
finiteness properties of \Cref{f:thermodynamics}; assume also, as there,
that \(H(\,\cdot\,,d)\) is concave on \((0,\infty)\) for every \(d>0\).
If
\[
 (s_1,\rho),(s_2,\rho)\in\partial Q(\beta_c,h_c),
 \qquad \rho>0,\quad s_1<s_2,
\]
then \(H(\,\cdot\,,\rho)\) has finite one-sided derivatives at \(\beta_c\)
and
\[
 H'_+(\beta_c,\rho)\le-s_2<-s_1\le H'_-(\beta_c,\rho).
\]
\end{lemma}

\begin{proof}
The field part of either support inequality says that \(\rho\) supports
\(Q(\beta_c,\cdot)\) at \(h_c\).
Density duality therefore gives
\(H(\beta_c,\rho)=h_c\rho-Q(\beta_c,h_c)\).
The full support inequality, with \(s=s_i\), reads
\[
 Q(\beta,h)\ge Q(\beta_c,h_c)
                 +s(\beta-\beta_c)+\rho(h-h_c).
\]
Rearrange and take the supremum over \(h\).
The reverse duality gives
\[
 H(\beta,\rho)\le H(\beta_c,\rho)-s(\beta-\beta_c).
\]
The function \(H(\,\cdot\,,\rho)\) is finite and concave on
\((0,\infty)\), so its one-sided derivatives at an interior point
exist and are finite.
Dividing the last inequality by positive and negative increments yields
the stated ordering.
\end{proof}

\begin{proof}[Proof of \Cref{thm:main}]
Take the potential \eqref{eq:final-potential} and the constants in
\Cref{m:tail}.
Write \(g=(g_1,\rho_0)\), where \(\rho_0=5p/3>0\), and set
\[
 \eta=\min\{\gamma/2,\rho_0/2\},\qquad
 I=(\rho_0-\eta,\rho_0+\eta)\subset(0,\infty).
\]
For every \(\rho\in I\), the two points
\[
 (g_1-\gamma/2,\rho),\qquad(g_1+\gamma/2,\rho)
\]
lie in \(\clB{g}{\gamma}\), since their squared distances from \(g\)
are \((\gamma/2)^2+(\rho-\rho_0)^2<\gamma^2/2\).
By \eqref{eq:limiting-support-ball}, both support the same pressure at
the same \((\beta_c,h_c)\), independently of \(\rho\).
Lemma~\ref{m:canonical-corner} therefore gives finite one-sided
canonical slopes whose difference is at least \(\gamma\) for every
\(\rho\in I\).
The all-sequence canonical limit of \Cref{f:thermodynamics} applies
in particular to \(N_L=\lfloor\rho L^3\rfloor\) for every \(\beta>0\).
The specified free energy equals
\[
 f(\beta,\rho)=\frac{t}{\beta}H_\Phi(\beta,\rho).
\]
For each \(\rho\in I\), it is finite for every \(\beta>0\).
The finite concave function \(H_\Phi(\,\cdot\,,\rho)\) is continuous
at \(\beta_c\); hence the common derivative term from \(1/\beta\)
cancels in the difference, giving
\[
 f'_-(\beta_c,\rho)-f'_+(\beta_c,\rho)
 =\frac{t}{\beta_c}
       \bigl(H'_{\Phi,-}(\beta_c,\rho)
                    -H'_{\Phi,+}(\beta_c,\rho)\bigr)
 \ge\frac{t\gamma}{\beta_c}>0.
\]
The potential and the density interval are fixed before the limiting
critical parameter is selected, so this one \(\beta_c\) works throughout
\(I\).
\end{proof}

